\section{共形块与 WZW：这个对应是什么级别}
\label{sec:blocks}

这是全篇最容易被说过头的地方。用户会问：Witten 到底证明了“三维 Chern--Simons 理论 = 二维 WZW 模型”，还是只证明了“两边的态空间维数一样”？答案在两者之间，而且必须把“对象”“函子”“范畴”三个层级分开。

\subsection{两条独立的路线}

WZW 一侧：一个二维共形场论，由 Wess--Zumino 项加 WZW 项定义，其手征代数（仿射李代数）的表示给出共形块。BPZ 方程、Tsuchiya--Kanie 的工作、Segal 的模函子，是这一侧的数学语言 \src{352}（文献 [11][16][32][39]）。

CS 一侧：第 \ref{sec:canonical} 节把三维理论沿 $\Sigma$ 切开，得到 $\cH_\Sigma^{(J)}=H^0(\cM_J,\cL^{\otimes k})$。

问题是：这两边是不是同一个东西？

\subsection{原文说了什么}

Witten 在 §3.1 说这些态空间“与 Segal 的模函子是同一物” \src{366}；在 §3.2 直接把 $\cH_\Sigma$ 与 BPZ 的共形块空间等同 \src{370}。英文用的是 “\textbf{are}”，不是 “isomorphic to”——从措辞看，他认为是同一个对象，不是碰巧同构。

但证据只到\textbf{对象级}：两边都等于 $H^0(\cM_J,\cL^{\otimes k})$，带同一个射影平坦丛 \src{369--370}。他没有证明函子性（缝合相容），§5 只有一句 “It is obvious at this point” \src{397}。

原文最强的\textbf{数学}断言在 §5 \src{396}：按 Segal 的看法，把这个辛结构量子化 $LG/G$，得到的正是圈群的基本不可约最高权表示。但这一句归功于 Segal（文献 [39]），而且限于 genus 0 与基本表示。

\subsection{三个层级}

\begin{center}
\small
\begin{tabular}{@{}p{2.3cm}p{5.2cm}p{6.2cm}@{}}
\toprule
层级 & 断言 & 原文证据 \\
\midrule
对象级 & 两边的态空间是同一个向量空间 & p.370 “are”；两边都是 $H^0(\cM_J,\cL^{\otimes k})$ \\
函子级 & 复结构变化时，两边给出同一个射影平坦丛；缝合相容 & p.370 只给丛与联络；缝合只到 p.397 “It is obvious” \\
范畴级 & 两个理论作为 TQFT/CFT 等价 & 未证；framing、边界条件、配分函数层面都有余量 \\
\bottomrule
\end{tabular}
\end{center}

所以正确的说法是：\textbf{强于“维数相同”，弱于“范畴等价”}。写“Witten 证明了两个理论等价”是错的；写“只是维数巧合”也错。Moore 中译的措辞更精确：向量空间同构，在 K\"ahler 量子化时是模空间上全纯向量丛的等同 \src{Moore 中译 p.76}。

\subsection{level 的账：裸 $k$ 与有效 $k+h^\vee$}

WZW 侧的一切公式里出现的是 $k+h^\vee$，不是裸 $k$。来源不是“两个理论能级不同”，而是 \S2 的单圈修正 \src{359}：(2.16)(2.17) 里 $\eta$ 的度规依赖把耦合平移为 $k+c_2(G)/2$，而 $c_2(SU(N))=2N$，故平移量是 $h^\vee=N$。Witten 自己说，这“似乎解释了下述事实：1+1 维流代数的许多公式中，量子修正的效果正是把 $k$ 替换为 $k+c_2(G)/2$” \src{359}。

在几何量子化的语言里，同一个平移来自半形式修正（第 \ref{sec:canonical} 节末）。最终 $SU(N)$ 的量子参数是 \eqref{eq:q}，$SU(2)$ 是 $q=e^{2\pi\ii/(k+2)}$ \src{382}。

\begin{keypoint}
\textbf{术语建议}：把 $k$ 叫“裸能级”、$k+h^\vee$ 叫“有效能级”。本地译本有一处写“WZW 能级是 $k+N$”，按这个口径是“有效能级”；若严格按模 $S$ 矩阵的 $q$ 定义“能级”，则 $k+h^\vee$ 才配称能级。两种说法不冲突，但混用会与可积性判据 $2j\le k$ 打架（后者用的是裸 $k$）。
\end{keypoint}

\subsection{KZ 方程：正文没有，但就在旁边}

KZ（Knizhnik--Zamolodchikov）方程在本文正文里\textbf{零出现}；最接近的是 p.370 引 BPZ 的微分方程 \src{370}。但教材里 KZ 正是 §3.2 那条射影平坦联络在 $g=0$、$n$ 个标记点时的显式形式：共形块丛 $\mathcal{E}_\lambda\to\mathrm{Conf}_n(\mathbb{CP}^1)$ 上的联络 $\nabla=\dd-\sum_i L^{(i)}_{-1}\dd z_i$，水平截面满足 KZ 方程，系数是 $\frac{1}{k+2}\frac{\Omega^{(ij)}}{z_i-z_j}$ \src{河野 p.25--28}。

所以本文与 KZ 的关系是：本文给“存在一条平坦联络”，KZ 给“这条联络长什么样”。不要把后者读进本文。

\begin{boundary}
\begin{itemize}[leftmargin=1.4em]
\item 原文给出：对象级等同（p.370）、射影平坦丛（p.370）、Segal 定理的引用（p.396）、三点维数 $N_{ijk}$（p.373--374）。
\item 我们补的：三层级划分；“裸能级/有效能级”的术语建议；KZ 与本文关系的定位。
\item 未证：函子性（p.397 只有 “obvious”）；范畴级等价；半形式修正的符号（$K^{1/2}$ 给 $k-h^\vee$ 还是 $k+h^\vee$，取决于 $\omega$ 与 $\mathrm{Det}$ 的取向约定，原文未给）。
\end{itemize}
\end{boundary}
